\problem{6.18}{Why translation and parity generate free-particle degeneracy}
\textbf{(a)} For any test function $f$,
\[
(\PP\T(a)f)(x)=f(-x-a),\qquad
(\T(a)\PP f)(x)=f(a-x).
\]
These are not generally equal; instead
\begin{equation}
\boxed{\PP\T(a)=\T(-a)\PP.}
\end{equation}
\textbf{(b)} The momentum eigenfunction
$f_p(x)=\ee^{\ii px/\hb}/\sqrt{2\pi\hb}$
obeys
\[
\PP f_p(x)=f_p(-x)=f_{-p}(x).
\]
The free Hamiltonian $H=\p^2/(2m)$
commutes with parity. Thus if
$Hf_p=E_pf_p$,
$H(\PP f_p)=\PP Hf_p=E_p f_{-p}$.
Hence
\begin{equation}
\boxed{E_{-p}=E_p=p^2/(2m),\qquad p\ne0.}
\end{equation}
\textbf{(c)} Define the even and odd
continuum-normalized functions
$c_p(x)=\cos(px/\hb)/\sqrt{\pi\hb}$,
$s_p(x)=\sin(px/\hb)/\sqrt{\pi\hb}$.
Angle-addition formulae give
\[
\begin{pmatrix}\T(a)c_p\\\T(a)s_p\end{pmatrix}
=
\begin{pmatrix}
\cos(pa/\hb)&\sin(pa/\hb)\\
-\sin(pa/\hb)&\cos(pa/\hb)
\end{pmatrix}
\begin{pmatrix}c_p\\s_p\end{pmatrix}.
\]
They are mixed by translations, so
they must belong to the same energy
eigenspace if translations commute
with $H$.
\phys{The momentum and parity bases
diagonalize different commuting
pairs; their ability to mix at fixed
energy reveals the degeneracy.}

\problem{6.19}{Commutators of a vector ladder operator}
Let $V_\pm=V_x\pm\ii V_y$,
and assume vector commutators
$[L_i,V_j]=\ii\hb\epsilon_{ijk}V_k$.
\textbf{(a)} Direct substitution yields
\[
[L_z,V_x]=\ii\hb V_y,\qquad
[L_z,V_y]=-\ii\hb V_x,
\]
so
\begin{equation}
\boxed{[L_z,V_\pm]=\pm\hb V_\pm.}
\end{equation}
Use $L^2=L_-L_++L_z^2+\hb L_z$
and
\[
[L_+,V_+]=0,\quad
[L_-,V_+]=-2\hb V_z,\quad
[L_z,V_+]=\hb V_+.
\]
Expanding $[AB,C]=A[B,C]+[A,C]B$
and commuting $L_z$ through the
vector components gives
\begin{equation}
\boxed{[L^2,V_\pm]=
2\hb^2V_\pm
\pm2\hb V_\pm L_z
\mp2\hb V_zL_\pm.}
\end{equation}
\textbf{(b)} Suppose
$L^2\ket{\ell,\ell}
=\hb^2\ell(\ell+1)\ket{\ell,\ell}$.
Since $L_+\ket{\ell,\ell}=0$,
the previous formula gives
\[
L^2 V_+\ket{\ell,\ell}
=\hb^2[\ell(\ell+1)+2+2\ell]
V_+\ket{\ell,\ell}.
\]
The factor is $(\ell+1)(\ell+2)$.
Also
$L_zV_+\ket{\ell,\ell}
=(\ell+1)\hb V_+\ket{\ell,\ell}$.
Therefore either
$V_+\ket{\ell,\ell}=0$, or it has
\begin{equation}
\boxed{\ell'=\ell+1,\qquad
m'=\ell+1.}
\end{equation}
This is the highest-weight special
case of the rank-one tensor
selection rules.

\problem{6.20}{Scalar selection rule using the lowering operator}
A scalar $f$ has $[L_-,f]=0$.
Fix $\ell'=\ell$ and consider
$\ket{n\ell m}$ and $\ket{n'\ell,m-1}$.
The matrix element of the commutator is
\[
0=\bra{n'\ell,m-1}[L_-,f]
\ket{n\ell m}.
\]
For a scalar, $[L_z,f]=0$, so only
equal magnetic indices occur.
Using
$L_-\ket{\ell m}=
\hb B^\ell_m\ket{\ell,m-1}$,
$B^\ell_m=\sqrt{(\ell+m)(\ell-m+1)}$,
we obtain
\[
0=\hb B^\ell_m
\left(
\bra{n'\ell m}f\ket{n\ell m}
-\bra{n'\ell,m-1}f\ket{n\ell,m-1}
\right).
\]
For $m>-\ell$ the prefactor is
nonzero, giving
\begin{equation}
\boxed{\bra{n'\ell m}f\ket{n\ell m}
=\bra{n'\ell,m-1}f\ket{n\ell,m-1}.}
\end{equation}
Thus the diagonal-in-$m$ matrix element
is independent of $m$, the same
selection relation obtained from
$[L_+,f]=0$.

\problem{6.21}{A scalar reduced matrix element for a hydrogen superposition}
Take the normalized
$\ket\psi=(\ket{211}+\ket{21,-1})/\sqrt2$.
Rotational scalar selection rules
for the magnitude $r$ state
\[
\bra{n'\ell'm'}r\ket{n\ell m}
=\delta_{\ell'\ell}\delta_{m'm}
\mathcal R_{n'\ell;n\ell},
\]
where $\mathcal R_{n'\ell;n\ell}$
is the single reduced radial
matrix element for fixed $\ell$.
The cross terms have $m'\ne m$ and
vanish. The diagonal ones are equal,
so
\[
\av r=\frac12(\mathcal R_{21;21}
+\mathcal R_{21;21})=\mathcal R_{21;21}.
\]
For hydrogen's $2p$ radial function
$R_{21}(r)=
r\ee^{-r/(2a_0)}/(2\sqrt6\,a_0^{5/2})$,
\begin{align}
\mathcal R_{21;21}
&=\int_0^\infty
r^3|R_{21}(r)|^2\dd r\\
&=\frac1{24a_0^5}
\int_0^\infty r^5\ee^{-r/a_0}\dd r
=\frac{5!a_0^6}{24a_0^5}.
\end{align}
Consequently
\begin{equation}
\boxed{\av r=5a_0.}
\end{equation}
\phys{Rotational symmetry reduces two
apparently separate diagonal
expectation values to one radial integral.}

\problem{6.22}{Vector-operator commutators and magnetic selection rule}
\textbf{(a)} Start from
$[L_i,V_j]=\ii\hb\epsilon_{ijk}V_k$.
Using $L_\pm=L_x\pm\ii L_y$
and $V_\pm=V_x\pm\ii V_y$,
expand the commutators componentwise.
The results are
\begin{align}
\boxed{[L_z,V_z]}&=0,&
\boxed{[L_z,V_\pm]}&=\pm\hb V_\pm,\\
\boxed{[L_\pm,V_\pm]}&=0,&
\boxed{[L_\pm,V_z]}&=\mp\hb V_\pm,\\
\boxed{[L_\pm,V_\mp]}&=\pm2\hb V_z.
\end{align}
For example
$[L_+,V_z]=[L_x,V_z]
+\ii[L_y,V_z]=-\ii\hb V_y-\hb V_x
=-\hb V_+$; the remaining identities
follow analogously.
\textbf{(b)} Insert eigenstates
$\ket a=\ket{n\ell m}$ and
$\ket b=\ket{n'\ell'm'}$
into $[L_z,V_z]=0$:
\[
0=\bra b[L_z,V_z]\ket a
=\hb(m'-m)\bra b V_z\ket a.
\]
Therefore
\begin{equation}
\boxed{\bra{n'\ell'm'}V_z\ket{n\ell m}
=0\quad\text{unless }m'=m.}
\end{equation}
Likewise $V_+$ needs $m'=m+1$,
and $V_-$ needs $m'=m-1$.
These are the magnetic quantum-number
selection rules in Equations 6.56--6.58.

\problem{6.23}{Clebsch--Gordan coefficient recursions}
\textbf{(a)} In the product space
$\ket{j_1j_2m_1m_2}$,
\[
\ket{JM}=\sum_{m_1+m_2=M}
C_{m_1m_2;M}^{j_1j_2;J}
\ket{j_1j_2m_1m_2}.
\]
Take the inner product with a fixed
uncoupled basis vector.
Orthonormality leaves exactly one term:
\begin{equation}
\boxed{C_{m_1m_2;M}^{j_1j_2;J}
=\braket{j_1j_2m_1m_2}{JM}.}
\end{equation}
\textbf{(b)} Define ladder coefficients
\[
A_m^j=\sqrt{(j-m)(j+m+1)},\qquad
B_m^j=\sqrt{(j+m)(j-m+1)}.
\]
Apply $J_+=J_{1+}+J_{2+}$ to the
coupled decomposition of $\ket{JM}$.
The left-hand side is
$\hb A_M^J\ket{J,M+1}$.
Match the coefficient of the
uncoupled $\ket{m_1,m_2}$:
\begin{equation}
\boxed{
A_M^J C^{J,M+1}_{m_1m_2}=
B_{m_1}^{j_1}C^{J,M}_{m_1-1,m_2}
+B_{m_2}^{j_2}C^{J,M}_{m_1,m_2-1}.}
\end{equation}
Similarly, $J_-$ yields
\begin{equation}
\boxed{
B_M^J C^{J,M-1}_{m_1m_2}=
A_{m_1}^{j_1}C^{J,M}_{m_1+1,m_2}
+A_{m_2}^{j_2}C^{J,M}_{m_1,m_2+1}.}
\end{equation}
In these formulas the fixed superscripts
$j_1,j_2,J$ on the $C$ coefficients
are suppressed for clarity.
Out-of-range magnetic indices have
zero coefficient. This establishes
both recursions directly from angular
momentum addition.

\problem{6.24}{Six vector matrix-element recursions and Wigner--Eckart}
Define
\[
M^{q}_{m'm}=\bra{n'\ell'm'}V_q\ket{n\ell m},
\quad q=+,-,z,
\]
and $A_m^\ell,B_m^\ell$
as in Problem 6.23.
\textbf{(a)} Sandwiching the six
commutators $[L_\pm,V_\pm]=0$,
$[L_\pm,V_z]=\mp\hb V_\pm$,
$[L_\pm,V_\mp]=\pm2\hb V_z$
between the angular eigenstates
produces the six identities
\begin{align}
B_{m'}^{\ell'}M^+_{m'-1,m}
-A_m^\ell M^+_{m',m+1}&=0,\\
A_{m'}^{\ell'}M^-_{m'+1,m}
-B_m^\ell M^-_{m',m-1}&=0,\\
B_{m'}^{\ell'}M^z_{m'-1,m}
-A_m^\ell M^z_{m',m+1}&=-M^+_{m'm},\\
A_{m'}^{\ell'}M^z_{m'+1,m}
-B_m^\ell M^z_{m',m-1}&=+M^-_{m'm},\\
B_{m'}^{\ell'}M^-_{m'-1,m}
-A_m^\ell M^-_{m',m+1}&=+2M^z_{m'm},\\
A_{m'}^{\ell'}M^+_{m'+1,m}
-B_m^\ell M^+_{m',m-1}&=-2M^z_{m'm}.
\end{align}
The terms on the left are obtained by
moving the ladder operator to the bra
or acting on the ket; the sign on the
right comes from its commutator.
\textbf{(b)} Adopt the spherical tensor
convention $V_{\pm1}=\mp V_\pm/\sqrt2$,
$V_0=V_z$. The Wigner--Eckart
ansatz (with phase convention of the text) is
\begin{equation}
\boxed{\begin{aligned}
M^+_{m'm}&=-\sqrt2\,
C^{\ell',m'}_{\ell,m;1,1}\,\mathcal R,\\
M^-_{m'm}&=+\sqrt2\,
C^{\ell',m'}_{\ell,m;1,-1}\,\mathcal R,\\
M^z_{m'm}&=
C^{\ell',m'}_{\ell,m;1,0}\,\mathcal R,
\end{aligned}}
\end{equation}
where $\mathcal R=\bra{n'\ell'}V\ket{n\ell}$
is a reduced matrix element independent of
$m,m'$ and tensor component.
Insert these three expressions into
the six commutator identities.
In each equation the common factor
$\mathcal R$ cancels, while the remaining
Clebsch--Gordan terms are precisely
the raising or lowering recursions
derived in Problem 6.23.
Hence the Wigner--Eckart form
solves all six constraints.
\phys{Rotations fix every dependence
on $m,m'$; radial dynamics is carried
only by the reduced matrix element.}

\problem{6.25}{The dipole of a $2p$ hydrogen superposition}
Take
\[
\ket\psi=(\ket{211}+\ket{210})/\sqrt2,\qquad
\mathbf p_e=-e\mathbf r.
\]
Each orbital has $\ell=1$, so
$\PP\ket\psi=-\ket\psi$;
the entire superposition has definite
odd parity. Because
$\PP^\dagger\mathbf p_e\PP=-\mathbf p_e$,
\[
\bra\psi\mathbf p_e\ket\psi
=\bra\psi\PP^\dagger\mathbf p_e\PP\ket\psi
=-\bra\psi\mathbf p_e\ket\psi.
\]
Therefore
\begin{equation}
\boxed{\av{p_{e,x}}=\av{p_{e,y}}=
\av{p_{e,z}}=0.}
\end{equation}
In Wigner--Eckart notation every
component is proportional to
the same reduced dipole matrix element
$\bra{2,1}\mathbf p_e\ket{2,1}$,
which is \emph{zero by Laporte's rule}.
Thus the vanishing is independent
of the Clebsch--Gordan numerical values;
no radial integration is required.
